\section{Introduction}

The goals of this analysis are investigation of hadron production, properties of QGP, QGP phase space, heavy-ion and light collision system evolution.
While measurements in this analysis by themselves might not provide a new valuable insight into any of aforementioned goals, they are most useful as complementary measurements of previous results for different observables, collision systems, energies, and particles species.
$K^{*0}(892)$ and $\bar{K}^{*0}(892)$ mesons were chosen as probes. 
The decay mode of interest for this analysis is $\bar{K^{*0}}(892) \rightarrow \pi^\pm K^\mp$ and $\bar{K^{*0}}(892) \rightarrow \pi^\pm K^\mp$ with $\text{BR} \approx 0.66601$ (\href{https://pdg.lbl.gov/2025/web/viewer.html?file=../listings/rpp2025-list-K-star-892.pdf}{PDG 2025}). 
Since $K^{*0}(892)$ and $\bar{K}^{*0}(892)$ are experimentally indistinguishable in this text it will be simply called \Kstar from now on. Results of this analysis are invariant \pt spectra and nuclear modification factors \rab and \rcp of \Kstar as functions of \pt in 5 centrality classes: 0-20\%, 20-40\%, 40-60\%, 60-88\%, 0-88\%. 
\pt ranges of yields measurements were chosen to be the same as in \Kstar production study in $p+p$ (\href{https://phenix-intra.sdcc.bnl.gov/phenix/WWW/p/draft/pshukla/AnaKstarNew/anaKstar5.pdf}{AN911}) so that calculation of \rab would be possible without any altering of measured \Kstar production results in $p+p$.
